\documentclass[tikz,border=2pt]{standalone}
\usepackage{amsmath,amssymb}
\usepackage{pgfplots}
\pgfplotsset{compat=1.18}

\begin{document}
% The tangent function is a bijection from the bounded interval (-pi/2, pi/2) onto all of R,
% so a short interval already has the cardinality of the continuum.
\begin{tikzpicture}
	\begin{axis}[
		axis lines=middle,
		xmin=-2.3, xmax=2.4, ymin=-3.4, ymax=3.6,
		xtick=\empty, ytick=\empty,
		xlabel={$x$}, ylabel={$y$},
		xlabel style={below right}, ylabel style={above left},
		samples=200, smooth, clip=true,
		width=8cm, height=6cm
		]
		\draw[gray, dashed] (axis cs:-1.5708,-3.4) -- (axis cs:-1.5708,3.6);
		\draw[gray, dashed] (axis cs:1.5708,-3.4) -- (axis cs:1.5708,3.6);
		\addplot[thick, blue, domain=-1.28:1.28] {tan(deg(x))};
		\node[blue, anchor=east, font=\small] at (axis cs:-0.12,1.3) {$y=\tan x$};
		\node[anchor=north east, font=\small] at (axis cs:-1.6,-0.08) {$-\tfrac\pi2$};
		\node[anchor=north west, font=\small] at (axis cs:1.6,-0.08) {$\tfrac\pi2$};
	\end{axis}
\end{tikzpicture}
\end{document}
